\documentclass[11pt]{article}
\usepackage[margin=1in]{geometry}
\usepackage{amsmath,amssymb,amsthm}
\usepackage{hyperref}
\usepackage{enumitem}

\title{SAT Rapidity Frontier: Current Formal State and Remaining Bridge}
\author{HQIV Lean notes}
\date{\today}

\begin{document}
\maketitle

\section*{Purpose}

This note summarizes the current formal state of the SAT rapidity program in the
repository, from the bottom-level certified facts already in Lean up to the final
missing geometric bridge. The goal is to phrase the setup clearly enough to think
through the last step.

\section{Bottom layer: what is already certified}

The SAT rapidity stack is organized as \emph{several} layers (worst-case semantics,
shared manifold / rapidity, residual control, and later---packing, direction selection,
exact-count/plane bridge, ribbon-cover collapse).  The first three are summarized here.

\subsection{SAT worst-case certification layer}

In \texttt{Hqiv/Geometry/SATWorstCaseCertified.lean}, the repository already has:

\begin{itemize}[leftmargin=2em]
  \item sound-prune semantics: pruning never removes a satisfying assignment;
  \item survivor-exhaustion semantics: exhaustive evaluation of the surviving pool
  with no model implies UNSAT;
  \item work-envelope transfer: if survivor work is bounded by a baseline plus an
  additive residual budget $\varepsilon$, and $\varepsilon$ is within the usual
  root-scale budget, then the normalized survivor work lies inside the same
  \(1 + n^{1/n}\) envelope used on the ATSP side;
  \item cumulative arity-residual transfer:
  \[
    \sum_i \varepsilon_i \leq \text{baselineWork} \cdot n^{1/n}
    \quad \Longrightarrow \quad
    \frac{\text{survivorWork}}{\text{baselineWork}} \leq 1 + n^{1/n};
  \]
  \item a polynomial-budget hook:
  if both baseline work and cumulative residual budget are polynomially bounded,
  then survivor work is polynomially bounded.
\end{itemize}

So the SAT file already gives the \emph{formal landing zone} for a polynomial-time-style
argument. What it does \emph{not} supply is the geometric theorem that forces those
polynomial bounds.

\subsection{Shared manifold / rapidity layer}

In \texttt{Hqiv/Geometry/SATRapidityManifold.lean}, we now formalize the idea that:

\begin{itemize}[leftmargin=2em]
  \item variables and clauses both embed into a shared ambient manifold;
  \item there is one common scalar rapidity observable on that manifold;
  \item the variable-side and clause-side rapidity profiles are both read from that
  same observable;
  \item a common rapidity threshold can therefore control both sides at once.
\end{itemize}

At the current scaffold level, rapidity is phrased in the simplest successor-step
form: a discrete \(m \mapsto m+1\) combinatorial step read from the shared manifold.

Concretely, the file now contains:

\begin{itemize}[leftmargin=2em]
  \item a shared manifold structure \texttt{SATSharedManifold};
  \item rapidity profiles on the variable and clause sides;
  \item successor-step rapidity channels \texttt{variableRapidityStep} and
  \texttt{clauseRapidityStep};
  \item common-threshold lemmas in both profile and successor-step form;
  \item successor-style envelope constants and transfer lemmas for
  \(\text{varDim}+1\) and \(\text{clauseDim}+1\), mirroring the ATSP \(n+1\) route.
\end{itemize}

\subsection{Residual-control layer}

The same file also introduces the first serious theorem-start for the missing bridge:

\begin{quote}
\textbf{Successor-step residual control:} each SAT residual is indexed by either a
variable-side or clause-side successor step, and that residual is bounded by the
corresponding rapidity step, which is itself bounded by one common threshold.
\end{quote}

Formally this is the structure \texttt{SuccessorStepResidualControl}. From it, the
repository already proves:

\begin{enumerate}[leftmargin=2em]
  \item every residual is bounded by the common threshold;
  \item therefore one can build a \texttt{SATSharedRapidityCertificate};
  \item if the threshold is polynomially bounded and the residual list length is
  polynomially bounded, then the cumulative residual budget is polynomially bounded.
\end{enumerate}

Thus rapidity can \emph{state} successor-step control cleanly; what a complete SAT
instance theorem still needs is a \emph{geometric or encoding witness} linking that
control to a bounded frontier (the combinatorial implications of such a witness are
largely already packaged higher in the stack---see \S\ref{sec:constructed}).

\section{Why stars-and-bars is probably not enough}

If the frontier bound is treated as a generic occupancy count, one gets something like
stars-and-bars behavior. That is too weak for the intended result: it does not force a
sharp collapse of the frontier, and it does not naturally explain why a common rapidity
threshold should drastically reduce the candidate family.

So the current working view is:

\begin{quote}
The missing theorem is probably not a bare combinatorial counting lemma. It is more
likely a packing theorem.
\end{quote}

In other words, the residual directions should be thought of as a family of directions,
rays, tangent vectors, or shell points on a common manifold, and the threshold should
force them to be sufficiently separated that only finitely many can coexist.

\section{Frontier geometry: the new packing layer}

\subsection{Analytical arc, annular ribbon, and moiré}

On the analytical sphere, the relevant continuum between a pole and a \(k\)-arity pole
is still a \textbf{one-dimensional arc} \(\gamma\). In Lean, the current counting kernel
is expressed through a local 2-dimensional section
(\texttt{Hqiv/Geometry/SATRapidityAnnulusCircle.lean}); this should be read as an
\emph{osculating ribbon section}, not as a global claim that the ambient manifold itself
is literally planar or globally spherical. In that local section, the arc lies on the target
shell circle \(C_M\). Discrete candidates lie in a tubular \textbf{ribbon} (points
within distance \(\varepsilon\) of some \(\gamma(t)\)); \textbf{moiré} here means
that a lattice point in the ribbon generically does \emph{not} lie on \(\gamma\)
itself. The radial annulus \(A_M(\tau)\) captures these near-arc points once
\(\varepsilon \le \tau\); this inclusion is proved in Lean from the triangle
inequality (\texttt{inAnnulus\_of\_shell\_arc\_ribbon}).

The current bridge file \texttt{Hqiv/Geometry/SATRapidityPlaneBridge.lean} makes this
local nature explicit: \texttt{PlaneWitnessMap} is the data \texttt{ResidualPoint M → Plane},
and the same object is also named \texttt{LocalRibbonSectionWitness} (definitionally the
same structure, with identity coercions) to stress that the global manifold remains abstract
while the counting kernel lives in a chosen osculating \(2\)-plane / local ribbon section.

In \texttt{Hqiv/Geometry/SATRapidityPackingBridge.lean}, the next frontier has now been
formalized.

The new object is a packing-style certificate:

\begin{quote}
\texttt{SATRapidityPackingCertificate}
\end{quote}

with the intended interpretation:

\begin{itemize}[leftmargin=2em]
  \item the residual frontier lives in an effective dimension \(d\);
  \item there is a packing / kissing-number style bound \(\mathrm{frontierBound}(d)\);
  \item the residual list length is bounded by that frontier bound;
  \item that bound is controlled by the combined variable/clause dimension.
\end{itemize}

The main theorem already proved from that scaffold is:

\begin{quote}
packing control + successor-step residual control
\(\Longrightarrow\) polynomial residual budget.
\end{quote}

So, structurally, the chain is now:

\[
\text{shared manifold}
\Longrightarrow
\text{successor-step rapidity}
\Longrightarrow
\text{residual domination}
\Longrightarrow
\text{packing-controlled frontier length}
\Longrightarrow
\text{polynomial residual budget}
\Longrightarrow
\text{SAT work certification.}
\]

\noindent
The same file stack also supports a \emph{parallel} route through planar exact-count bounds
(\(K_{\mathrm{exact}} \le 2|Q|\)) and \texttt{RibbonCoverCollapseData}; that route does not require
an abstract kissing-number function \(K(d)\) once the local shell hypotheses hold.

\section{What is now fully constructed in Lean}
\label{sec:constructed}

Since the original version of this note, the intermediate bridge has become substantially
more explicit.

\subsection{Direction selection, gap bridge, and geometric collapse are packaged}

The repository now contains:

\begin{itemize}[leftmargin=2em]
  \item \texttt{DirectionSelectionCertificate}: residuals choose canonical shell directions,
  with injectivity/separation packaged as explicit fields;
  \item \texttt{SATRapidityGapBridge}: the single record representing the statement
  ``rapidity gap controls canonical directions'';
  \item \texttt{SATRapidityGeometricCollapse}: the end-stage record packaging
  logarithmic effective dimension, sphere-code control, and successor-step bookkeeping.
\end{itemize}

Moreover, once the appropriate hypotheses are supplied (including via the plane-cover and
ribbon-cover constructors in \texttt{SATRapidityPlaneBridge.lean}), one can build all three levels
by explicit constructors rather than only by isolated \texttt{theorem} calls.

\subsection{Exact-count route is now explicit}

The exact-count route through local shell intersections has been formalized in
\texttt{SATRapidityExactCardinality.lean} and wired to the planar model in
\texttt{SATRapidityAnnulusCircle.lean} / \texttt{SATRapidityPlaneBridge.lean}.

The key proved ingredients are:

\begin{enumerate}[leftmargin=2em]
  \item for each nonzero shell center \(q\), a local shell fiber satisfying the
  \texttt{planeLocalShellIntersections} predicate obeys \(|I|\le 2\) (circle--circle geometry);
  \item therefore the exact admissible-direction count
  \[
    K_{\mathrm{exact}} := \#\Bigl(\bigcup_{q \in Q} \mathrm{intersections}(q)\Bigr)
  \]
  satisfies \(K_{\mathrm{exact}} \le 2 |Q|\);
  \item if residual list length is dominated by \(K_{\mathrm{exact}}\), then it is bounded by
  \(2|Q|\), and hence by \(2B\) whenever \(|Q| \le B\);
  \item when \(Q\) is the carrier of an \texttt{ArcRibbonLatticeCardBound} aligned with the
  annulus model, the same chain yields \(\mathrm{len} \le 2\cdot\mathrm{countBound}(\cdot)\)
  without choosing an intermediate \(B\) by hand;
  \item natural-number and real-cast versions feed the \texttt{hSphereLen} slot once combined
  with a frontier comparison into \texttt{sphereCodeBound effectiveDim}.
\end{enumerate}

The bridge file contains \emph{several} explicit constructor families, all hypothesis-explicit:

\begin{itemize}[leftmargin=2em]
  \item \textbf{Comparison frontier:} \texttt{DirectionSelectionCertificate.of\_plane\_cover\_bound},
  \texttt{SATRapidityGapBridge.of\_plane\_cover\_bound},
  \texttt{SATRapidityGeometricCollapse.of\_plane\_cover\_bound}, assuming
  \(2B \le \mathrm{sphereCodeBound}(\mathrm{effectiveDim})\) (as reals).
  \item \textbf{Exact frontier:} the parallel \texttt{of\_plane\_cover\_exact\_frontier} constructors
  when \(\mathrm{sphereCodeBound}(\mathrm{effectiveDim}) = 2B\), so no separate inequality step
  is needed at that dimension.
  \item \textbf{Threading only:}
  \texttt{direction\_selection\_with\_plane\_witness\_implies\_gap\_bridge} and
  \texttt{direction\_selection\_with\_plane\_witness\_implies\_geometric\_collapse}
  repackage an \emph{existing} certificate and only carry a
  \texttt{PlaneWitnessMap} for joint hypotheses (they do not derive \texttt{hSphereLen} from the plane).
\end{itemize}

A small finite-size lemma internal to the same file records that the factorization/oracle
\texttt{candidateList} has length exactly \(3(\texttt{candidateStepBound}(n)+1)\)
(\texttt{factorization\_candidate\_family\_card\_eq}), packaging the list-length part of that
pipeline into the SAT rapidity stack (membership in an annulus family is still separate).

\subsection{Single-hypothesis collapse of steps 2--5}

The strongest consolidated packaging is the structure

\begin{quote}
\texttt{RibbonCoverCollapseData M D control}
\end{quote}

in \texttt{SATRapidityPlaneBridge.lean}.  It assumes an ambient
\texttt{DirectionSelectionCertificate} \texttt{D} and \texttt{SuccessorStepResidualControl}
\texttt{control}, and bundles:

\begin{itemize}[leftmargin=2em]
  \item \texttt{witness} : \texttt{LocalRibbonSectionWitness M} (same as \texttt{PlaneWitnessMap});
  \item \texttt{family} : \texttt{AnnulusLatticeFamily} at \((\texttt{shellRadius},\texttt{thresholdWidth})\) from \texttt{D.annulusModel};
  \item \texttt{intersections} : \texttt{Plane → Finset Plane} and the fiber hypotheses
  \texttt{hne}, \texttt{hloc} (\texttt{planeLocalShellIntersections} on each center);
  \item cover \texttt{hCover} : residual centers land in \texttt{family.carrier};
  \item exact-count domination \texttt{hK} : \(\texttt{D.residuals.length} \le K_{\mathrm{exact}}\);
  \item polynomial slot \texttt{countBound}, \texttt{varDim}, \texttt{clauseDim}, \texttt{tauBits}, and
  \texttt{hCard} : \(|\texttt{family.carrier}| \le \texttt{countBound}(\cdots)\);
  \item real comparison \texttt{hCountFrontier} :
  \(2\cdot\texttt{countBound}(\cdots) \le \texttt{D.sphereCodeBound}(\texttt{effectiveDim})\);
  \item successor packaging \texttt{hShared}, \texttt{hLen}, and \texttt{hPoly} for the sphere-code step.
\end{itemize}

From this record, Lean defines:

\begin{itemize}[leftmargin=2em]
  \item \texttt{RibbonCoverCollapseData.hSphereLen} (theorem): discharges the sphere-length inequality for \texttt{D};
  \item \texttt{toDirectionSelectionCertificate}: \texttt{D} with \texttt{hSphereLen} replaced by that proof;
  \item \texttt{toGapBridge}, \texttt{toGeometricCollapse};
  \item named entry points \texttt{ribbon\_cover\_collapse},
  \texttt{ribbon\_cover\_collapse\_to\_gap\_bridge},
  \texttt{ribbon\_cover\_collapse\_to\_direction\_selection},
  \texttt{ribbon\_cover\_collapse\_to\_sphere\_code};
  \item \texttt{ribbon\_cover\_collapse\_hasPolynomialResidualBudget}: closes into
  \texttt{HasPolynomialResidualBudget} under the usual nonnegativity / threshold hypotheses;
  \item \texttt{ribbon\_cover\_collapse\_implies\_nat\_root\_envelope}: combines the collapsed residual budget
  with the certified ATSP-style envelope (\texttt{ATSPWorstCaseCertified}) for a single numeric worst-case bound.
\end{itemize}

The method \texttt{toGeometricCollapse} applies \texttt{DirectionSelectionCertificate.toGeometricCollapse} to the
ambient \texttt{D}; the method \texttt{toDirectionSelectionCertificate} replaces \texttt{D.hSphereLen} with the
theorem \texttt{RibbonCoverCollapseData.hSphereLen}.  In applications where the ribbon witness is the
authoritative justification for the sphere-length bound, the upgraded certificate should be used when
passing into collapse or sphere-code packaging.

So the downstream pipeline after \texttt{RibbonCoverCollapseData} is fully wired in Lean; the
remaining substantive task is still to \emph{construct} that witness (and the scaffold interfaces
below) from first-principles geometry and encoding.

\section{What the final missing piece must say}

At this point, the \emph{formal} pipeline after a suitable witness is largely implemented
in Lean (\S\ref{sec:constructed}).  What remains unproved is sharp in a different sense:
one must \emph{construct} the witness from SAT geometry and encoding, not reprove the
downstream certificate algebra.

One seeks a theorem of the following shape.

\subsection*{Desired frontier theorem}

There exists a natural effective dimension \(d\) attached to the shared variable/clause
manifold such that the residual directions generated by the SAT oracle admit a local
ribbon-section cover, and hence form a packing/counting family on a shell, tangent sphere,
or comparable local geometric carrier. The number of admissible residual directions is then
bounded by a function \(K(d)\) of packing/kissing-number or exact-shell-count type.

If additionally \(K(d)\) is polynomially controlled by the combined input dimension
\(\text{varDim}+\text{clauseDim}\), then the entire residual frontier is polynomially bounded.

\subsection*{Why this is the right target}

This would supply what the formal pipeline is already set up to consume:

\begin{enumerate}[leftmargin=2em]
  \item identify each residual with a local ribbon-section center and associated shell data;
  \item show thresholded rapidity forces the hypotheses used in the planar/exact-count layer
  (nontrivial centers, \texttt{planeLocalShellIntersections} fibers, cover in a finite annulus family);
  \item establish the exact-count domination \(\mathrm{len} \le K_{\mathrm{exact}}\) and a polynomial
  cardinality bound on \(\#Q\) (e.g.\ via \texttt{ArcRibbonLatticeCardBound});
  \item package these as a \texttt{RibbonCoverCollapseData} witness (or equivalent inputs to
  \texttt{of\_plane\_cover\_bound} / \texttt{of\_plane\_cover\_exact\_frontier});
  \item then \emph{already formalized}: \texttt{RibbonCoverCollapseData} \(\to\)
  \texttt{SATRapidityGeometricCollapse} \(\to\) \texttt{SATRapiditySphereCodeCertificate} \(\to\)
  \texttt{HasPolynomialResidualBudget} (see \texttt{ribbon\_cover\_collapse\_hasPolynomialResidualBudget}),
  and hence the certified SAT work / polynomial-budget story at the worst-case layer.
\end{enumerate}

\noindent
A separate, older abstraction \texttt{SATRapidityPackingCertificate} in
\texttt{SATRapidityPackingBridge.lean} still represents a generic packing/kissing-style frontier
bound; the ribbon-cover route in \texttt{SATRapidityPlaneBridge.lean} closes through the
sphere-code certificate instead.  Either style of frontier input can feed polynomial-budget
machinery once the numeric hypotheses match the downstream lemmas.

\medskip
\noindent
Equivalently: the \emph{implications} after a ribbon-cover witness are already proved; what is
still missing is the \emph{theorem that produces} that witness from first principles.

\section{Interpretive note on dimension and collapse}

The motivating speculation is that the bridge may depend on a nontrivial dimension regime,
not monotone growth with dimension. In that picture, one does \emph{not} want a naive
``more dimensions means more space'' argument. Rather, one wants a regime in which the
shared manifold geometry imposes a sharp structural bound on admissible frontier directions.

That is why sphere packing / kissing-number heuristics are more promising here than generic
stars-and-bars counting. The right theorem would exploit geometry strong enough to turn the
frontier into a constrained packing problem instead of an unconstrained occupancy problem.

\section{Two concrete candidate approaches for the final bridge}

The current best two candidate approaches are the following.

\subsection{Angular separation on the rapidity sphere / local section}

This is the most direct geometric route.

\begin{enumerate}[leftmargin=2em]
  \item Map each residual to a unit vector or local shell direction whose direction is the gradient of the
  shared rapidity observable (or its local ribbon-section representative).
  \item Use the common threshold to force angular separation between distinct
  residual directions. A model inequality would be
  \[
    |\cos \theta_{ij}| \leq 1 - \delta
  \]
  for distinct residuals \(i \neq j\).
  \item View the residual family as a spherical code on the unit sphere in
  \(\mathbb{R}^d\), where \(d\) is derived from
  \(\mathrm{varDim}+\mathrm{clauseDim}\).
  \item Apply a spherical-code or kissing-number bound \(K(d)\) to conclude that
  the frontier cardinality is at most \(K(d)\).
\end{enumerate}

This is attractive because it connects almost immediately to the existing packing
certificate abstraction: \(K(d)\) can serve as the \texttt{frontierBound} in
\texttt{SATRapidityPackingCertificate}.

\subsection{Riemannian tangent-space packing}

This is the more general manifold-native route.

\begin{enumerate}[leftmargin=2em]
  \item Equip the shared manifold with a Riemannian metric structure compatible
  with the existing oracle language.
  \item Interpret residuals as tangent vectors with norm bounded by the successor-step
  rapidity.
  \item Use the threshold to enforce a minimum angle between any two such tangent
  vectors.
  \item Apply a tangent-space packing lemma: for example, one has exponential-type
  coarse bounds such as \(2^d\) under large pairwise angular separation, and then one
  tries to improve these to polynomial-in-\(d\) bounds under additional curvature,
  injectivity-radius, or shell-regularity assumptions.
\end{enumerate}

This route is more flexible, because it can absorb curvature and local manifold data,
but the final theorem may need stronger geometric hypotheses than the spherical-code route.

\section{What is still genuinely missing}

Despite the progress above, a fully internal proof still lacks the following substantive ingredient.

\subsection{The one remaining witness-producing theorem}

The current Lean development still does \emph{not} derive from first principles that:

\begin{enumerate}[leftmargin=2em]
  \item every residual lies in a concrete local ribbon-section annulus family,
  \item the associated shell fibers are the right admissible local intersections,
  \item residual length is dominated by the resulting exact-count union model, and
  \item the resulting annulus family carries the required polynomial cardinality witness.
\end{enumerate}

This whole cluster is isolated as the construction of one \texttt{RibbonCoverCollapseData}
witness (or, equivalently, the inputs to the \texttt{of\_plane\_cover\_bound} constructor family
together with the separate hypotheses that define a ribbon-cover collapse in the formal sense).

\subsection{What is still scaffold-level}

Two definitions remain intentionally lightweight placeholders in the abstract SAT scaffold:

\begin{itemize}[leftmargin=2em]
  \item \texttt{localShellIntersections} on the fully abstract manifold side is still the empty-finset placeholder;
  \item \texttt{angle} is still the constant-\(0\) placeholder on abstract shell directions.
\end{itemize}

These do \emph{not} break the current bridge construction, because the new exact-count work is done in
the local ribbon-section model. But a fully geometric final theorem should eventually replace those
placeholder interfaces by real shell-intersection and angle data, or else prove that the local ribbon-section
construction is the only geometry actually needed.

\subsection*{How they fit the current Lean scaffold}

A full internal proof still has to supply a frontier bound and the encoding-side hypotheses above.
Once those are available, the \textbf{primary wired endpoint} in the current development is:

\begin{itemize}[leftmargin=2em]
  \item build a \texttt{RibbonCoverCollapseData} witness (or the equivalent plane-cover inputs);
  \item obtain \texttt{SATRapidityGeometricCollapse} and \texttt{SATRapiditySphereCodeCertificate};
  \item apply \texttt{ribbon\_cover\_collapse\_hasPolynomialResidualBudget} to reach
  \texttt{HasPolynomialResidualBudget}.
\end{itemize}

\noindent
Separately, a classical packing-style certificate \texttt{SATRapidityPackingCertificate} can still
serve as an abstract \texttt{frontierBound} route into polynomial-budget lemmas when its numeric
hypotheses are arranged to match \texttt{packingCertificate\_hasPolynomialResidualBudget}.

\medskip
\noindent
So the open choice is less about ``which Lean record to use at the end'' than about \emph{which
geometric theorem produces the frontier and cover data}.  A spherical-code bound, a tangent-space
packing bound, or the proved planar exact-count \(\le 2\#Q\) chain are different \emph{sources} for
those hypotheses; the repository already implements the bookkeeping once the bounds are supplied.

\section{Best anchor points for the next proof attempt}

After reviewing the agent-facing documentation, the best anchor points for the remaining work are
now fairly clear.

\subsection{Primary anchor: manifold rapidity roadmaps}

The strongest direct anchor is the manifold roadmap:

\begin{itemize}[leftmargin=2em]
  \item \texttt{AGENTS/MANIFOLD\_ZETA\_ROADMAP.md}
\end{itemize}

This document already states, in essentially the right form, that the missing theorem is a
\emph{ray theorem}: rapidity should factor through or agree with a canonical map from shell index
to a unit tangent direction on a Riemannian slice. That is almost exactly the geometric-collapse
problem now represented by \texttt{SATRapidityGeometricCollapse}.

So the most natural theorem target is not a totally new idea, but a SAT-specialized version of the
same roadmap question:

\begin{quote}
Can the shared SAT rapidity observable be shown to factor through a low-dimensional family of
canonical directions (rays / tangent vectors / shell points) on the manifold, with enough angular
separation to force a packing bound?
\end{quote}

\subsection{Secondary anchor: theorem index for existing reusable pieces}

The next best anchor is:

\begin{itemize}[leftmargin=2em]
  \item \texttt{AGENTS/THEOREMS.md}
\end{itemize}

The most relevant existing names are:

\begin{itemize}[leftmargin=2em]
  \item Euclidean and shell-volume anchors from \texttt{SpatialSliceManifold};
  \item the rapidity / polar-angle identification from
  \texttt{RapidityZetaPhaseBridge};
  \item the explicit shell-to-\(m+1\) monotonicity and shell-surface proxies from
  \texttt{OctonionSphereConstruction} and related modules;
  \item the Fourier patch / moir\'e slope language listed in the archive and modular-curvature notes.
\end{itemize}

These suggest that the next proof should be phrased less as ``find a new object'' and more as
``bind together already-existing shell, phase, and patch concentration channels into one direction-selection theorem.''

\subsection{Third anchor: modular / Fourier patch narrative}

The file

\begin{itemize}[leftmargin=2em]
  \item \texttt{AGENTS/MODULAR\_THETA\_CURVATURE\_BRIDGE.md}
\end{itemize}

is especially relevant because it explicitly points to Fourier patch concentration and curved-manifold
lifts via heat kernels / Laplacian spectrum. This is a strong hint that, if the SAT rapidity frontier
really collapses, a plausible mechanism may be:

\begin{enumerate}[leftmargin=2em]
  \item moir\'e / patch concentration isolates a low-dimensional family of dominant directions,
  \item rapidity phase selects among those directions,
  \item packing then bounds how many such dominant directions can coexist.
\end{enumerate}

\subsection{Fourth anchor: Euler/Gauss-style shell intersection counts}

There is also a simpler and potentially very powerful anchor: lattice intersection counts per shell.

The existing corpus does not yet contain a full Gauss circle-problem theorem in dimension two, but it
does contain the right style of ingredients:

\begin{itemize}[leftmargin=2em]
  \item discrete shell counts such as \texttt{r8} in
  \texttt{Hqiv/Algebra/IntegerLatticeShellCount8.lean};
  \item shell monotonicity and continuous sphere/ball proxies in
  \texttt{OctonionSphereConstruction};
  \item Euclidean shell / slice geometry in \texttt{SpatialSliceManifold} and
  \texttt{EuclideanBallHorizontalSlice};
  \item the explicit warning in the docs that current combinatorics are incremental shell counts,
  not yet full Gauss-circle totals.
\end{itemize}

This suggests another framing of the direction-selection theorem:

\begin{quote}
Instead of proving the frontier bound directly from abstract packing, prove that the annulus-shell
construction reduces the residual frontier to a controlled lattice intersection problem on one shell,
and then bound the number of admissible intersections by an Euler/Gauss-style shell-count function.
\end{quote}

In the local 2D picture, each residual point in the annulus probes the target shell by a unit circle,
and the canonical directions are the shell intersection points. The natural counting question is then:

\begin{quote}
How many such intersection points can occur on a fixed shell under the threshold and separation rules?
\end{quote}

That is extremely close in spirit to a circle-problem or shell-count question, and may give a simpler
route than a fully general manifold packing theorem if the relevant shell counts can be controlled sharply.

\subsection*{Recommended next Lean move}

The next substantive Lean goal is unchanged in \emph{name}---produce a
\texttt{RibbonCoverCollapseData} witness (or equivalent \texttt{of\_plane\_cover\_*} inputs) from
shared-manifold, rapidity, and encoding hypotheses---but the \emph{downstream} implications are already
implemented (\S\ref{sec:constructed}), so new code should focus on geometry and \texttt{hK}, not on
re-proving certificate projections.

If the Euler/Gauss-shell route is pursued, a companion module such as
\texttt{Hqiv/Geometry/SATRapidityShellIntersectionCount.lean} could package shell-intersection counting as an
alternate source of frontier bounds feeding the same constructors.

\section{Current status in one sentence}

The repository has the full formal scaffold from shared manifold rapidity through direction selection,
gap bridge, and geometric collapse; the exact-count / planar-fiber inequalities and the
\texttt{RibbonCoverCollapseData} packaging (including closure into \texttt{HasPolynomialResidualBudget}
and the ATSP-style root envelope) are implemented in Lean. The remaining substantive work is to
construct the ribbon-cover / encoding witness and to replace abstract-placeholder \texttt{localShellIntersections}
and \texttt{angle} with geometric content where a global statement is desired.

\section*{Files to consult}

\begin{itemize}[leftmargin=2em]
  \item \texttt{Hqiv/Geometry/SATWorstCaseCertified.lean}
  \item \texttt{Hqiv/Geometry/SATRapidityManifold.lean}
  \item \texttt{Hqiv/Geometry/SATRapidityDirectionSelection.lean}
  \item \texttt{Hqiv/Geometry/SATRapidityExactCardinality.lean}
  \item \texttt{Hqiv/Geometry/SATRapidityAnnulusCircle.lean}
  \item \texttt{Hqiv/Geometry/SATRapidityPackingBridge.lean}
  \item \texttt{Hqiv/Geometry/SATRapiditySpherePacking.lean}
  \item \texttt{Hqiv/Geometry/SATRapidityTangentPacking.lean}
  \item \texttt{Hqiv/Geometry/SATRapidityGapBridge.lean}
  \item \texttt{Hqiv/Geometry/SATRapidityPlaneBridge.lean} (ribbon-cover collapse, constructors, polynomial closure)
  \item \texttt{Hqiv/Geometry/ATSPWorstCaseCertified.lean} (envelope used by \texttt{ribbon\_cover\_collapse\_implies\_nat\_root\_envelope})
  \item \texttt{Hqiv/Geometry/SpatialSliceRapidityScaffold.lean}
  \item \texttt{Hqiv/Geometry/GeneralRiemannianRapidityOracle.lean}
\end{itemize}

\end{document}